\subsection{Introduction \quad / \quad \textthai{บทนำ}}
\begin{paracol}{2}
Deep Learning (DL)\index{Deep Learning@ดีปเลิร์นนิง (Deep Learning)} has transitioned from a purely computational interest to a core methodology in contemporary physics. The fundamental premise of DL in scientific research is the \textbf{Universal Function Approximation Theorem}\index{Universal Function Approximation Theorem@ทฤษฎีการประมาณฟังก์ชันสากล (Universal Function Approximation Theorem)}, which posits that a neural network with a single hidden layer and a non-linear activation function can approximate any continuous function on a compact subset of $\mathbb{R}^n$. 

In the realm of physics, this is transformative. Traditional models often rely on explicit differential equations or \textit{Matched Filtering}\index{Matched Filtering@การกรองแบบจับคู่ (Matched Filtering)}, which requires a pre-defined template of the signal being sought. While powerful, matched filtering is computationally expensive and struggles with "non-stationary" noise—irregular fluctuations that do not follow a simple statistical distribution. 

This chapter focuses on a critical application: the detection of \textbf{Gravitational Waves}\index{Gravitational Waves@คลื่นความโน้มถ่วง (Gravitational Waves)}. First predicted by Einstein in 1916 and detected by the LIGO collaboration in 2015, these spacetime ripples are often millions of times fainter than the seismic noise surrounding them. We will demonstrate how 1D Convolutional Neural Networks (1D CNNs) can act as adaptive, non-linear filters to "see" through this noise, identifying binary black hole mergers with unprecedented speed.

\switchcolumn

\begin{thai}
ดีปเลิร์นนิง (Deep Learning) ได้เปลี่ยนจากความสนใจทางคอมพิวเตอร์อย่างหมดจด กลายเป็นระเบียบวิธีวิจัยหลักในฟิสิกส์ร่วมสมัย ข้อมติพื้นฐานของ DL ในการวิจัยทางวิทยาศาสตร์คือ \textbf{ทฤษฎีการประมาณฟังก์ชันสากล (Universal Function Approximation Theorem)} ซึ่งกล่าวว่าโครงข่ายประสาทเทียมที่มีเลเยอร์ซ่อนอยู่แม้เพียงชั้นเดียวและมีฟังก์ชันกระตุ้นแบบไม่เชิงเส้น สามารถประมาณค่าฟังก์ชันต่อเนื่องใดๆ บนเซตย่อยที่กระชับของ $\mathbb{R}^n$ ได้

ในทางฟิสิกส์ สิ่งนี้คือการเปลี่ยนแปลงเชิงโครงสร้าง โมเดลแบบดั้งเดิมมักอาศัยสมการเชิงอนุพันธ์ที่ชัดเจนหรือ \textit{\textbf{การกรองแบบจับคู่ (Matched Filtering)}\index{Matched Filtering@การกรองแบบจับคู่ (Matched Filtering)}} ซึ่งต้องการต้นแบบ (Template) ของสัญญาณที่กำหนดไว้ล่วงหน้า แม้จะมีประสิทธิภาพสูง แต่การกรองแบบจับคู่ต้องใช้ทรัพยากรการคำนวณมหาศาล และมีปัญหาเมื่อเจอกับสัญญาณรบกวนแบบ "ไม่คงที่" (Non-stationary noise)

บทนี้จะเน้นไปที่การประยุกต์ใช้ที่สำคัญยิ่ง: \textbf{คลื่นความโน้มถ่วง (Gravitational Waves)}\index{Gravitational Waves@คลื่นความโน้มถ่วง (Gravitational Waves)} ซึ่งไอน์สไตน์ทำนายไว้ในปี 1916 และตรวจพบโดยกลุ่มความร่วมมือ LIGO ในปี 2015 ระลอกคลื่นในอวกาศกาลเหล่านี้มักจะเบากว่าสัญญาณรบกวนจากแผ่นดินไหวรอบตัวพวกมันหลายล้านเท่า เราจะแสดงให้เห็นว่าโครงข่ายประสาทเทียมแบบคอนโวลูชัน 1 มิติ (1D CNNs) สามารถทำหน้าที่เป็นตัวกรองแบบปรับตัวและไม่เชิงเส้น เพื่อ "มองทะลุ" สัญญาณรบกวนเหล่านี้ และระบุการรวมตัวของหลุมดำคู่ด้วยความเร็วที่ไม่เคยมีมาก่อน
\end{thai}
\end{paracol}
